\documentclass[12pt,a4paper]{report}

\usepackage{fontspec}
\usepackage{polyglossia}
\setdefaultlanguage{romanian}
\usepackage[margin=2.25cm]{geometry}
\usepackage{setspace}
\usepackage{amsmath}
\usepackage{amssymb}
\usepackage{enumitem}
\usepackage{longtable}
\usepackage{array}
\usepackage{booktabs}
\usepackage{xcolor}
\usepackage{listings}
\usepackage[
    colorlinks,
    linkcolor={black},
    citecolor={black},
    urlcolor={blue}
]{hyperref}

\onehalfspacing
\setlength{\emergencystretch}{3em}
\setlength{\parskip}{0.15em}
\setlist{nosep,leftmargin=*}

\definecolor{codebackground}{rgb}{0.96,0.96,0.96}
\definecolor{codecomment}{rgb}{0.23,0.45,0.30}
\definecolor{codekeyword}{rgb}{0.10,0.20,0.55}
\definecolor{codestring}{rgb}{0.55,0.15,0.15}

\lstdefinestyle{codestyle}{
    backgroundcolor=\color{codebackground},
    basicstyle=\ttfamily\scriptsize,
    breaklines=true,
    frame=single,
    columns=fullflexible,
    keepspaces=true,
    showstringspaces=false,
    numbers=left,
    numberstyle=\tiny\color{gray},
    keywordstyle=\color{codekeyword}\bfseries,
    commentstyle=\color{codecomment},
    stringstyle=\color{codestring},
    tabsize=4
}
\lstset{style=codestyle}

\newcommand{\codechunk}[2]{%
\subsection*{Liniile #1--#2}
\lstinputlisting[language=Python,firstline=#1,lastline=#2,firstnumber=#1]{../src/models/mask_model.py}
}

\title{Explicarea fisierului \texttt{mask\_model.py}}
\author{}
\date{}

\begin{document}

\maketitle
\tableofcontents
\clearpage

\chapter{Rolul fisierului \texttt{mask\_model.py}}

Fisierul \path{src/models/mask_model.py} defineste arhitectura neuronala principala
a proiectului: MaskNet. Modelul este proiectat pentru imbunatatirea semnalului
de vorbire prin estimarea unei masti spectrale. Intrarea este de obicei o
spectrograma de magnitudine cu forma:

\[
x \in \mathbb{R}^{B \times C \times F \times T},
\]

unde \(B\) este batch size-ul, \(C\) numarul de canale, \(F\) numarul de bin-uri
de frecventa, iar \(T\) numarul de cadre temporale. Pentru modelul de magnitudine,
\(C=1\). Pentru varianta complexa, \(C=2\), deoarece intrarea contine partea
reala si partea imaginara a STFT.

Modelul foloseste o arhitectura de tip U-Net: encoderul reduce rezolutia
timp-frecventa si creste numarul de canale, bottleneck-ul extrage reprezentari
compacte, iar decoderul reconstruieste o masca la rezolutia initiala. Skip
connections pastreaza detalii locale, attention-ul poate repondera canalele, iar
BiLSTM-ul modeleaza dependente temporale.

\begin{longtable}{p{0.30\textwidth}p{0.60\textwidth}}
\toprule
\textbf{Clasa} & \textbf{Rol} \\
\midrule
\endhead
\texttt{ConvBlock} & Convolutie 2D, BatchNorm si ReLU. Bloc de baza pentru encoder
si decoder. \\
\texttt{DepthwiseSeparableConvBlock} & Varianta mai eficienta, separand convolutia
spatiala de proiectia pe canale. \\
\texttt{ResidualConvBlock} & Doua blocuri convolutionale cu conexiune reziduala. \\
\texttt{EncoderBlock} & Doua convolutii si optional downsampling prin MaxPool. \\
\texttt{DecoderBlock} & Upsampling, concatenare cu skip connection si convolutii. \\
\texttt{AttentionBlock} & Mecanism de channel attention. \\
\texttt{BiLSTMBlock} & Modeleaza dependente temporale in bottleneck. \\
\texttt{MaskNet} & Reteaua completa U-Net + optional LSTM + optional attention. \\
\texttt{SimpleMaskNet} & Varianta veche, pastrata pentru compatibilitate. \\
\bottomrule
\end{longtable}

\chapter{Importuri si blocul convolutional simplu}

\codechunk{1}{16}

\begin{description}
    \item[Liniile 1--3.] Importa PyTorch, modulul \texttt{torch.nn} si functiile
    neuronale din \texttt{torch.nn.functional}. \texttt{nn} este folosit pentru
    definirea layerelor, iar \texttt{F} pentru operatii functionale precum
    padding-ul sau ReLU in modelul simplu.

    \item[Liniile 6--7.] Defineste clasa \texttt{ConvBlock}, mostenind
    \texttt{nn.Module}. Acesta este blocul standard al arhitecturii: convolutie,
    normalizare si activare.

    \item[Liniile 8--13.] Constructorul primeste numarul de canale de intrare,
    canale de iesire, dimensiunea kernelului si stride-ul. Padding-ul este calculat
    ca \texttt{kernel\_size // 2}, ceea ce pastreaza dimensiunea spatiala cand
    stride-ul este 1. Blocul contine \texttt{Conv2d}, \texttt{BatchNorm2d} si
    \texttt{ReLU}.

    \item[Liniile 15--16.] \texttt{forward} aplica layer-ele in ordinea:
    convolutie, batch normalization, ReLU. Iesirea are activare neliniara si poate
    fi trimisa catre urmatorul bloc.
\end{description}

\chapter{Convolutie depthwise separabila}

\codechunk{19}{43}

\begin{description}
    \item[Liniile 19--20.] \texttt{DepthwiseSeparableConvBlock} este o alternativa
    mai eficienta la \texttt{ConvBlock}. Ideea este sa se separe filtrarea spatiala
    de combinarea canalelor.

    \item[Liniile 22--35.] Prima convolutie este \texttt{depthwise}: are
    \texttt{groups=in\_channels}, ceea ce inseamna ca fiecare canal este filtrat
    separat. Apoi \texttt{pointwise} foloseste kernel \(1 \times 1\) pentru a
    combina canalele si a produce \texttt{out\_channels}.

    \item[Liniile 36--37.] Sunt definite BatchNorm-urile si ReLU-ul comun. BatchNorm
    stabilizeaza distributia activarilor, iar ReLU introduce neliniaritate.

    \item[Liniile 39--43.] In \texttt{forward}, semnalul trece intai prin convolutia
    depthwise, apoi prin convolutia pointwise. Aceasta varianta poate reduce
    numarul de parametri si costul computational, util pentru rulari rapide sau
    real-time.
\end{description}

\chapter{Bloc rezidual}

\codechunk{46}{72}

\begin{description}
    \item[Liniile 46--47.] \texttt{ResidualConvBlock} defineste un bloc cu doua
    straturi convolutionale si o conexiune reziduala. Conexiunea reziduala ajuta
    propagarea gradientului si permite modelului sa invete corectii fata de intrare.

    \item[Liniile 48--59.] Constructorul alege tipul de bloc intern: depthwise sau
    clasic. Sunt create doua convolutii succesive si o activare ReLU.

    \item[Liniile 60--65.] Daca numarul canalelor de intrare difera de cel al
    canalelor de iesire, reziduul nu poate fi adunat direct. De aceea se creeaza
    \texttt{self.proj}, o proiectie \(1 \times 1\) urmata de BatchNorm, care
    aliniaza dimensiunile.

    \item[Liniile 67--72.] In \texttt{forward}, se salveaza reziduul, se aplica
    cele doua convolutii, apoi se aduna iesirea cu reziduul. ReLU-ul final produce
    activarea blocului.
\end{description}

\chapter{Encoderul U-Net}

\codechunk{75}{107}

\begin{description}
    \item[Liniile 75--76.] \texttt{EncoderBlock} este blocul de coborare al U-Net-ului.
    El extrage trasaturi si, optional, reduce rezolutia timp-frecventa.

    \item[Liniile 77--91.] Constructorul primeste canalele, flag-ul de downsampling
    si optiunile rezidual/depthwise. Daca \texttt{use\_residual=True}, foloseste
    un \texttt{ResidualConvBlock}; altfel foloseste doua blocuri convolutionale
    simple sau depthwise.

    \item[Liniile 92--94.] Daca downsampling-ul este activ, se creeaza un
    \texttt{MaxPool2d} cu kernel si stride 2. Aceasta injumatateste frecventa si
    timpul in reprezentarea intermediara.

    \item[Liniile 96--107.] \texttt{forward} produce doua iesiri: tensorul dupa
    pooling si tensorul \texttt{before\_pool}. Al doilea este skip connection-ul
    care va fi transmis decoderului pentru a recupera detalii fine.
\end{description}

\chapter{Decoderul U-Net}

\codechunk{110}{147}

\begin{description}
    \item[Liniile 110--111.] \texttt{DecoderBlock} este blocul de urcare al U-Net-ului.
    El mareste rezolutia si concateneaza informatia cu skip connection-ul din encoder.

    \item[Liniile 112--124.] Constructorul creeaza un \texttt{ConvTranspose2d}
    pentru upsampling. Apoi alege intre bloc rezidual si doua blocuri convolutionale
    standard. Dupa concatenare, numarul de canale devine dublu, deci primul bloc
    primeste \texttt{out\_channels * 2}.

    \item[Liniile 126--136.] \texttt{forward} aplica upsampling. Daca forma tensorului
    upsampled nu coincide perfect cu skip connection-ul, se calculeaza diferentele
    pe inaltime/frecventa si latime/timp, apoi se aplica \texttt{F.pad}. Aceasta
    trateaza diferentele cauzate de impartiri si rotunjiri la downsampling.

    \item[Liniile 138--147.] Tensorul upsampled este concatenat cu skip-ul pe
    dimensiunea canalelor. Apoi trece prin blocul rezidual sau prin doua convolutii.
    Rezultatul devine intrarea pentru urmatorul nivel de decoder.
\end{description}

\chapter{Channel Attention}

\codechunk{150}{168}

\begin{description}
    \item[Liniile 150--151.] \texttt{AttentionBlock} implementeaza un mecanism de
    attention pe canale, similar conceptual cu Squeeze-and-Excitation.

    \item[Liniile 152--161.] Constructorul creeaza un \texttt{AdaptiveAvgPool2d(1)},
    care reduce fiecare canal la o singura valoare medie. Apoi o retea fully
    connected mica reduce dimensiunea canalelor, aplica ReLU, revine la numarul
    original de canale si foloseste sigmoid pentru ponderi intre 0 si 1.

    \item[Liniile 163--168.] In \texttt{forward}, tensorul este redus la vectorul
    \((B, C)\), trecut prin reteaua FC si remodelat la \((B,C,1,1)\). Iesirea
    este produsul dintre tensorul original si ponderile de attention extinse.
    Astfel, canalele utile pot fi amplificate, iar cele mai putin relevante pot
    fi diminuate.
\end{description}

\chapter{Blocul BiLSTM}

\codechunk{171}{200}

\begin{description}
    \item[Liniile 171--172.] \texttt{BiLSTMBlock} adauga modelare temporala in
    bottleneck. Convolutiile observa vecinatati locale, iar LSTM-ul poate urmari
    dependente pe axa timpului.

    \item[Liniile 173--183.] Constructorul creeaza un LSTM bidirectional cu
    \texttt{batch\_first=True}. Dimensiunea de intrare este data de canalele si
    frecventele comprimate, iar dimensiunea ascunsa este \texttt{hidden\_size}.
    Pentru ca LSTM-ul este bidirectional, iesirea are dimensiune dubla si este
    proiectata inapoi prin \texttt{self.linear}.

    \item[Liniile 186--194.] In \texttt{forward}, tensorul are forma
    \((B,C,F,T)\). Este permutat la \((B,T,C,F)\), apoi aplatizat la
    \((B,T,C\cdot F)\), deoarece LSTM-ul trebuie sa primeasca o secventa temporala.

    \item[Liniile 196--200.] LSTM-ul proceseaza secventa, linearul readuce dimensiunea
    la cea originala, apoi tensorul este remodelat si permutat inapoi la forma
    \((B,C,F,T)\).
\end{description}

\chapter{Clasa MaskNet: constructorul}

\codechunk{204}{252}

\begin{description}
    \item[Liniile 204--219.] Docstring-ul descrie arhitectura MaskNet: U-Net cu
    skip connections, attention optional si BiLSTM optional. Argumentele principale
    controleaza canalele de intrare, dimensiunea bazei si componentele activate.

    \item[Liniile 220--236.] Constructorul primeste toate optiunile necesare:
    \texttt{in\_channels}, \texttt{base\_channels}, \texttt{use\_lstm},
    \texttt{use\_attention}, \texttt{use\_residual}, \texttt{use\_depthwise},
    numarul bin-urilor de frecventa, canalele de iesire si tipul mastii.

    \item[Liniile 237--248.] Salveaza configuratia ca atribute ale obiectului.
    Aceste atribute sunt folosite ulterior in \texttt{forward} si in activarea de
    iesire.

    \item[Liniile 250--252.] Alege clasa de bloc convolutional de baza:
    \texttt{DepthwiseSeparableConvBlock} daca \texttt{use\_depthwise=True}, altfel
    \texttt{ConvBlock}.
\end{description}

\codechunk{255}{303}

\begin{description}
    \item[Liniile 255--258.] Creeaza cele patru niveluri de encoder. Canalele cresc
    progresiv: \(C\), \(2C\), \(4C\), \(8C\), unde \(C=\texttt{base\_channels}\).
    Fiecare nivel reduce rezolutia prin pooling.

    \item[Liniile 261--271.] Creeaza bottleneck-ul. Daca modelul este rezidual,
    foloseste doua blocuri reziduale. Altfel, foloseste doua blocuri convolutionale.
    Numarul de canale urca pana la \(16C\).

    \item[Liniile 274--278.] Daca LSTM-ul este activ, se calculeaza dimensiunea
    frecventei dupa patru downsampling-uri. Pentru 257 bin-uri, impartirea intreaga
    la 16 da dimensiunea folosita la LSTM. Intrarea LSTM-ului este
    \(16C \cdot F_{\text{down}}\).

    \item[Liniile 281--286.] Daca attention-ul este activ, se creeaza cate un
    \texttt{AttentionBlock} pentru fiecare skip connection.

    \item[Liniile 289--292.] Creeaza cele patru niveluri de decoder. Acestea inverseaza
    procesul encoderului: \(16C \rightarrow 8C \rightarrow 4C \rightarrow 2C
    \rightarrow C\).

    \item[Liniile 295--303.] \texttt{out\_conv} transforma canalele finale in
    \texttt{output\_channels}. Ultima convolutie \(1 \times 1\) produce masca
    spectrala: un canal pentru masca de magnitudine sau doua canale pentru masca
    complexa.
\end{description}

\chapter{Activarea de iesire si forward-ul MaskNet}

\codechunk{305}{356}

\begin{description}
    \item[Liniile 305--312.] \texttt{\_apply\_output\_activation} aplica activarea
    finala. \texttt{sigmoid} limiteaza masca la intervalul \([0,1]\), potrivit
    pentru masti de magnitudine. \texttt{tanh} inmultit cu \texttt{output\_scale}
    este potrivit pentru masti complexe, unde valorile pot fi negative sau mai
    mari decat 1. Optiunile \texttt{identity}/\texttt{linear} lasa iesirea neschimbata.

    \item[Liniile 314--324.] Docstring-ul \texttt{forward} fixeaza forma intrarii:
    \((batch, channels, freq, time)\). Iesirea este o masca la aceeasi rezolutie
    timp-frecventa.

    \item[Liniile 326--329.] Intrarea trece prin cele patru niveluri de encoder.
    La fiecare nivel se pastreaza si skip connection-ul.

    \item[Liniile 332--335.] Tensorul cel mai comprimat trece prin bottleneck si,
    daca este activ, prin BiLSTM.

    \item[Liniile 338--342.] Daca attention-ul este activ, skip connection-urile
    sunt reponderate pe canale inainte de a fi folosite de decoder.

    \item[Liniile 344--347.] Decoderul reconstruieste rezolutia initiala, folosind
    skip connection-urile in ordine inversa.

    \item[Liniile 350--356.] \texttt{out\_conv} produce masca, iar activarea finala
    o transforma in forma potrivita tipului de model.
\end{description}

\chapter{SimpleMaskNet}

\codechunk{360}{383}

\begin{description}
    \item[Liniile 360--361.] \texttt{SimpleMaskNet} este o arhitectura mai veche,
    pastrata pentru compatibilitate. Ea nu are U-Net, skip connections, LSTM sau
    attention.

    \item[Liniile 362--372.] Constructorul defineste trei convolutii cu BatchNorm
    si o convolutie finala \(1 \times 1\), urmata de sigmoid.

    \item[Liniile 374--383.] \texttt{forward} aplica secvential cele trei blocuri
    convolutionale cu ReLU, produce masca prin \texttt{out\_conv} si o limiteaza
    cu sigmoid. Aceasta varianta este mai simpla si mai rapida, dar are capacitate
    expresiva mai mica decat MaskNet-ul complet.
\end{description}

\chapter{Concluzie}

\texttt{mask\_model.py} este fisierul care defineste nucleul neuronal al
proiectului. Configuratia din \texttt{config.py} decide ce varianta de model se
construieste, iar acest fisier contine clasele efective care transforma
spectrograma zgomotoasa intr-o masca spectrala.

Arhitectura MaskNet este flexibila: poate functiona ca model de magnitudine,
poate folosi masti complexe, poate activa LSTM pentru context temporal, attention
pentru reponderarea canalelor, blocuri reziduale pentru propagarea gradientului
si convolutii depthwise pentru eficienta. Aceasta flexibilitate permite proiectului
sa compare modele rapide cu modele mai expresive, pastrand aceeasi interfata de
training si inferenta.

\end{document}
